\documentclass[11pt]{article}

\usepackage[T1]{fontenc}
\usepackage[utf8]{inputenc}
\usepackage{lmodern}
\usepackage{microtype}
\usepackage[margin=1in]{geometry}
\usepackage{amsmath,amssymb,amsthm,mathtools}
\usepackage{booktabs,longtable,array}
\usepackage{enumitem}
\usepackage{xcolor}
\usepackage{listings}
\usepackage{url}
\usepackage{hyperref}
\usepackage[nameinlink,capitalise,noabbrev]{cleveref}

\hypersetup{
  colorlinks=true,
  linkcolor=blue!45!black,
  citecolor=green!35!black,
  urlcolor=blue!55!black,
  pdftitle={A Radix--Carry Continuation of the Alaoglu--Erdos Program},
  pdfauthor={OpenAI GPT-5.6 Pro},
  pdfsubject={Alaoglu--Erdos conjecture; radix reinterpretation; carry cocycles; p-adic dynamics}
}

\allowdisplaybreaks
\emergencystretch=3em
\setlength{\parindent}{1.3em}
\setlength{\parskip}{0.25em}
\setlist{itemsep=0.25em,topsep=0.35em}

\newtheorem{theorem}{Theorem}[section]
\newtheorem{proposition}[theorem]{Proposition}
\newtheorem{lemma}[theorem]{Lemma}
\newtheorem{corollary}[theorem]{Corollary}
\newtheorem{conjecture}[theorem]{Conjectural target}
\newtheorem{problem}[theorem]{Research problem}
\theoremstyle{definition}
\newtheorem{definition}[theorem]{Definition}
\newtheorem{example}[theorem]{Example}
\theoremstyle{remark}
\newtheorem{remark}[theorem]{Remark}
\newtheorem{warning}[theorem]{Warning}

\newcommand{\N}{\mathbb N}
\newcommand{\Nzero}{\mathbb N_0}
\newcommand{\Z}{\mathbb Z}
\newcommand{\Q}{\mathbb Q}
\newcommand{\R}{\mathbb R}
\newcommand{\Zp}{\mathbb Z_p}
\newcommand{\floor}[1]{\left\lfloor #1\right\rfloor}
\newcommand{\ceil}[1]{\left\lceil #1\right\rceil}
\newcommand{\fracpart}[1]{\left\{#1\right\}}
\newcommand{\vTwo}{v_2}
\newcommand{\vThree}{v_3}
\newcommand{\PhiMap}{\Phi}
\newcommand{\PsiMap}{\Psi}
\newcommand{\Rad}{\mathcal R}
\newcommand{\Cantor}{\mathcal C}
\newcommand{\carry}{\mathsf C}
\newcommand{\revcarry}{\mathsf K}
\newcommand{\thetaAE}{\vartheta}
\newcommand{\alphaAE}{\alpha}
\newcommand{\evidence}[1]{\textsf{[#1]}}
\newcommand{\baseline}{the unified report}

\DeclareMathOperator{\supp}{supp}
\DeclareMathOperator{\ord}{ord}
\DeclareMathOperator{\len}{len}
\DeclareMathOperator{\rad}{rad}

\definecolor{proofgreen}{RGB}{30,105,65}
\definecolor{targetorange}{RGB}{155,82,15}
\definecolor{codegray}{RGB}{247,247,247}

\lstset{
  basicstyle=\ttfamily\small,
  backgroundcolor=\color{codegray},
  frame=single,
  rulecolor=\color{black!20},
  columns=fullflexible,
  keepspaces=true,
  showstringspaces=false,
  breaklines=true,
  aboveskip=0.8em,
  belowskip=0.8em
}

\title{\textbf{A Radix--Carry Continuation of the Alaoglu--Erd\H{o}s Program}\\[0.45em]
\large Exact digit transducers, defect dynamics, and structural-prime constraints}
\author{OpenAI GPT-5.6 Pro}
\date{22 August 2026}

\begin{document}
\maketitle

\begin{abstract}
Let
\[
  \thetaAE=\log_2 3,\qquad
  \alphaAE=\frac{1}{\thetaAE}=\log_3 2.
\]
The Alaoglu--Erd\H{o}s conjecture is equivalent to the assertion that an
integer value of $M^{\thetaAE}$, with $M\in\N$, forces $M$ to be a power of
$2$.  The attached unified report already develops the algebraic,
transcendence-theoretic, determinant, finite-verification, primitive-generator,
and dynamical structure of this problem.  This continuation deliberately does
not repeat that material.  It develops a separate radix mechanism.

Write an integer in base $2$ and reinterpret the same finite digit word in
base $3$.  The resulting map $\PhiMap$ is an order-preserving embedding of
$\Nzero$ into the ternary Cantor integers.  More strongly, it is an exact
cross-prime isometry:
\[
  \vThree\bigl(\PhiMap(m)-\PhiMap(n)\bigr)=\vTwo(m-n)
  \quad(m\ne n).
\]
Its failures of additivity and multiplicativity are not errors but explicit
positive carry polynomials.  Reinterpreting canonical ternary digits in base
$2$ gives a reverse map $\PsiMap$ with negative carry polynomials.  These facts
lead to sharp Archimedean envelopes
\[
  \PhiMap(n)\le n^{\thetaAE},\qquad
  \PsiMap(n)\ge n^{\alphaAE},
\]
with equality only at powers of the source base.

For a hypothetical nonintegral solution
$M=2^x$, $A=3^x$, define
\[
  P_k=\PhiMap(M^k),\quad D_k=A^k-P_k,
  \qquad Q_k=\PsiMap(A^k),\quad E_k=Q_k-M^k.
\]
Then $D_k,E_k$ are positive integers for every $k\ge1$.  This report proves
exact nonlinear carry recurrences for them; an exact mantissa formula; the
uniform estimate
\[
  \frac{\fracpart{kx}}{3}
  \le \frac{D_k}{A^k}
  \le (\log 3)\fracpart{kx};
\]
a complete cross-prime valuation law for the Cantor shadow orbit $P_k$; and a
first-collision theorem for binary squaring which gives
\[
 \vThree\bigl(P_{2k}-P_k^2\bigr)
 =2k\vTwo(M)+\vTwo\bigl(U^k-1\bigr),
 \qquad M=2^{\vTwo(M)}U,\ U\text{ odd}.
\]
A nonlinear spectral-radius reformulation also emerges:
\[
  \lim_{k\to\infty}\PhiMap(M^k)^{1/k}=M^{\thetaAE};
\]
hence the original conjecture asks whether this digit-dynamical spectral
radius can be an integer only on the carry-free monomial orbits $M=2^n$.

These are rigorous paper proofs, accompanied by independent finite executable
checks.  They do not yet close the conjecture.  The remaining obstruction is
isolated as a simultaneous Archimedean/$3$-adic rigidity problem for the same
positive defect sequence, rather than another reformulation involving only
logarithms.
\end{abstract}

\begin{center}
\fbox{\begin{minipage}{0.92\textwidth}
\textbf{Status discipline.}
\evidence{P} means proved in this report by a complete mathematical argument.
\evidence{C} means independently checked by finite computation and is used
only as a sanity check.  \evidence{B} refers to an input already established
in the attached unified report.  \evidence{T} is a proposed closing lemma or
research target.  No claim announced by another model or research group is
used as evidence.  In particular, the Alaoglu--Erd\H{o}s conjecture remains
open at the end of this continuation.
\end{minipage}}
\end{center}

\tableofcontents

\section{Scope, interface to the baseline, and the new result map}

The unified report is already a large consolidation.  Repeating its survey of
primitive generators, common odd cores, algebraic auxiliary outputs,
determinant repulsion, higher differences, finite zero-free regions, and
classical transcendence consequences would make this document longer while
making the frontier less visible.  We therefore use only the following
minimal interface.

\begin{enumerate}[label=(B\arabic*)]
\item If $2^x$ and $3^x$ are positive integers and $x$ is not an integer,
then $x>0$ and $x$ is irrational.  In fact, standard transcendence theory
makes $x$ transcendental, but irrationality is enough for the dynamical
arguments below. \evidence{B}
\item We may put
\[
  M=2^x\in\N,\qquad A=3^x\in\N,
\]
so that
\[
  A=M^{\thetaAE},\qquad \thetaAE=\log_2 3.
\]
Conversely, an integer pair $M,A$ satisfying this relation gives a solution
$x=\log_2 M$. \evidence{B}
\item Integer translation acts by
$(M,A)\mapsto (2^rM,3^rA)$.  The baseline contains much sharper primitive
normalisations; below we use only this elementary compatibility. \evidence{B}
\end{enumerate}

Everything after this interface is developed independently.  The main new
objects and their evidence levels are as follows.

\begin{longtable}{>{\raggedright\arraybackslash}p{0.27\textwidth}
                  >{\raggedright\arraybackslash}p{0.57\textwidth}
                  >{\centering\arraybackslash}p{0.08\textwidth}}
\toprule
Object or theorem & Content & Status\\
\midrule
\endhead
Radix embedding $\PhiMap$ & Binary digits re-evaluated in base $3$; strict
order embedding and exact $2$-adic/$3$-adic isometry. & \evidence{P}\\
Carry cocycles & Exact positive correction polynomials for addition and
multiplication; exact negative corrections for $\PsiMap$. & \evidence{P}\\
Power envelopes & $\PhiMap(n)\le n^{\thetaAE}$ and
$\PsiMap(n)\ge n^{\alphaAE}$, with complete equality classification and
explicit gaps. & \evidence{P}\\
Counterexample defects & Positive integer sequences $D_k,E_k$ with exact
nonlinear recurrences and carry-budget inequalities. & \evidence{P}\\
Mantissa observable & Exact representation of $D_k/A^k$ by the fractional
rotation $\fracpart{kx}$ and a Cantor digit embedding; universal two-sided
bounds. & \evidence{P}\\
Valuation resonance & Pairwise $3$-adic distances of $P_k$ are exactly the
binary distances of $M^k$; outside an explicit LTE resonance set, valuations
of $D_k-D_\ell$ are forced. & \evidence{P}\\
Square first collision & The first nonzero ternary digit of the square-carry
polynomial occurs at a computable position and has coefficient $1$. &
\evidence{P}\\
Radix spectral radius & $\sup_k\PhiMap(M^k)^{1/k}=M^{\thetaAE}$, giving a
purely digit-dynamical equivalent form of the conjecture. & \evidence{P}\\
Finite executable audit & Exhaustive checks of the exact identities over
large finite boxes; never used in a proof. & \evidence{C}\\
Closing theorem & A cross-place rigidity theorem strong enough to make the
carry budget incompatible with the rotation gauge. & \evidence{T}\\
\bottomrule
\end{longtable}

\section{Canonical digits and radix reinterpretation}

\subsection{General definition}

For an integer base $b\ge2$ and $n\in\Nzero$, write the canonical finite
expansion
\[
  n=\sum_{i=0}^{L} d_i^{(b)}(n)b^i,
  \qquad 0\le d_i^{(b)}(n)<b.
\]
Leading zeroes are harmless, but the canonical expansion has
$d_L^{(b)}(n)\ne0$ when $n>0$.

\begin{definition}[Radix reinterpretation]
For bases $b,c\ge2$, define
\[
  \Rad_{b\to c}(n)
   :=\sum_{i\ge0}d_i^{(b)}(n)c^i.
\]
The two maps used throughout are
\[
  \PhiMap:=\Rad_{2\to3},\qquad
  \PsiMap:=\Rad_{3\to2}.
\]
Thus $\PhiMap$ replaces the place values $1,2,4,\ldots$ of a binary word by
$1,3,9,\ldots$, while $\PsiMap$ evaluates canonical ternary digits
$0,1,2$ at the place values $1,2,4,\ldots$.
\end{definition}

The image of $\PhiMap$ consists precisely of the nonnegative integers whose
ternary digits belong to $\{0,1\}$.  We call them \emph{ternary Cantor
integers}.  On this image, $\PsiMap$ is the genuine inverse:
\[
  \PsiMap(\PhiMap(n))=n.
\]
On arbitrary integers, however, $\PsiMap$ is not monotone and is not an
inverse in either direction.  For example,
\[
  8=(22)_3,\quad \PsiMap(8)=6,
  \qquad 9=(100)_3,\quad \PsiMap(9)=4.
\]
This warning prevents a tempting but invalid order argument later.

\subsection{Strict order preservation in the expanding direction}

\begin{proposition}[Order embedding]\label{prop:order}
If $2\le b<c$, then $\Rad_{b\to c}$ is strictly increasing on $\Nzero$.
In particular, $\PhiMap$ is a strict order embedding.
\end{proposition}

\begin{proof}
Let $m<n$, and let $r$ be the highest digit position at which their canonical
base-$b$ expansions differ.  At that position the digit of $n$ exceeds the
digit of $m$ by at least $1$.  Lower positions can oppose this difference by
at most
\[
  (b-1)\sum_{i=0}^{r-1}c^i
  =(b-1)\frac{c^r-1}{c-1}<c^r,
\]
because $b<c$.  Hence the contribution $c^r$ at the highest differing
position dominates every possible lower-digit loss, so
$\Rad_{b\to c}(n)>\Rad_{b\to c}(m)$.
\end{proof}

\subsection{The exact cross-prime isometry}

The strongest elementary feature of $\PhiMap$ is not its monotonicity but its
preservation of the index of the first differing digit.

\begin{theorem}[Binary--ternary valuation isometry]\label{thm:isometry}
For distinct $m,n\in\Nzero$,
\[
  \boxed{\quad
  \vThree\bigl(\PhiMap(m)-\PhiMap(n)\bigr)
  =\vTwo(m-n).
  \quad}
\]
Equivalently, the first binary digit at which $m$ and $n$ differ has exactly
the same index as the first ternary digit at which their Cantor images differ.
\end{theorem}

\begin{proof}
Put $r=\vTwo(m-n)$.  Congruence modulo $2^r$ says that the binary digits in
positions $0,\ldots,r-1$ agree.  Failure of congruence modulo $2^{r+1}$ says
that the digits in position $r$ differ, necessarily by $\pm1$.  Therefore
\[
 \PhiMap(m)-\PhiMap(n)
 =3^r\bigl(\pm1+3z\bigr)
\]
for some $z\in\Z$.  The parenthesised factor is a $3$-adic unit, proving the
claim.
\end{proof}

\begin{corollary}[A $2$-adic/$3$-adic Cantor embedding]\label{cor:padic-extension}
The digit formula extends continuously to
\[
  \widehat\Phi:\Z_2\longrightarrow\Z_3,
  \qquad
  \widehat\Phi\!\left(\sum_{i\ge0}\varepsilon_i2^i\right)
  =\sum_{i\ge0}\varepsilon_i3^i,
  \quad \varepsilon_i\in\{0,1\}.
\]
It is a homeomorphism from $\Z_2$ onto the closed ternary Cantor subset of
$\Z_3$, and for $u\ne v$,
\[
  \vThree\bigl(\widehat\Phi(u)-\widehat\Phi(v)\bigr)=\vTwo(u-v).
\]
\end{corollary}

\begin{proof}
The series converges in $\Z_3$.  Agreement of the first $r$ binary digits is
equivalent to agreement modulo $2^r$ and becomes agreement of the first $r$
ternary digits, hence agreement modulo $3^r$.  The argument of
\cref{thm:isometry} gives equality at the first differing digit.  Compactness
of $\Z_2$ and the explicit inverse on $\{0,1\}$-digit strings give the
homeomorphism statement.
\end{proof}

This isometry is the first genuinely cross-place input of the report.  It
turns a binary multiplicative orbit into a ternary Cantor shadow with exactly
prescribed $3$-adic geometry.

\section{Carry cocycles: exact failure of ring homomorphism}

The map $\PhiMap$ is neither additive nor multiplicative.  Its defects are,
however, completely explicit and have a fixed sign on nonnegative integers.

\subsection{Addition}

Let $u_i,v_i,w_i\in\{0,1\}$ be the binary digits of $u,v,u+v$.  Let
$\kappa_{-1}=0$ and let $\kappa_i$ be the carry leaving position $i$, so
\[
  u_i+v_i+\kappa_{i-1}=w_i+2\kappa_i.
\]

\begin{theorem}[Addition carry identity]\label{thm:add-carry}
For $u,v\in\Nzero$,
\[
 \boxed{
 \PhiMap(u+v)=\PhiMap(u)+\PhiMap(v)
       +\sum_{i\ge0}\kappa_i3^i.}
\]
Consequently $\PhiMap$ is superadditive, with equality exactly when the binary
addition is carry-free.
\end{theorem}

\begin{proof}
Multiply the digit recurrence by $3^i$ and sum over $i$:
\[
 \PhiMap(u)+\PhiMap(v)+\sum_{i\ge0}\kappa_{i-1}3^i
 =\PhiMap(u+v)+2\sum_{i\ge0}\kappa_i3^i.
\]
The shifted carry sum is $3\sum_i\kappa_i3^i$, giving the formula after
rearrangement.  Every coefficient is nonnegative, and the correction vanishes
exactly when every carry vanishes.
\end{proof}

\subsection{Multiplication}

Let $r_i=\sum_{p+q=i}u_pv_q$ be the raw binary convolution coefficient.  On
normalising the convolution to binary digits $w_i$, carries satisfy
\[
  r_i+\kappa_{i-1}=w_i+2\kappa_i.
\]

\begin{theorem}[Multiplication carry identity]\label{thm:mul-carry}
For $u,v\in\Nzero$,
\[
 \boxed{
 \PhiMap(uv)=\PhiMap(u)\PhiMap(v)
       +\sum_{i\ge0}\kappa_i3^i.}
\]
Thus $\PhiMap$ is supermultiplicative.  Equality holds precisely when the raw
binary convolution already has every coefficient in $\{0,1\}$.
\end{theorem}

\begin{proof}
The same summation as in \cref{thm:add-carry} applies, with
$\sum_i r_i3^i=\PhiMap(u)\PhiMap(v)$.
\end{proof}

\begin{definition}[Carry cocycles]
Set
\[
 \mathfrak a(u,v)=\PhiMap(u+v)-\PhiMap(u)-\PhiMap(v),
 \qquad
 \mathfrak m(u,v)=\PhiMap(uv)-\PhiMap(u)\PhiMap(v).
\]
On $\Nzero^2$ these are nonnegative integers whose ternary values are the
carry polynomials in \cref{thm:add-carry,thm:mul-carry}.
\end{definition}

The terminology ``cocycle'' is useful: associativity forces identities among
these corrections.  For example,
\[
 \mathfrak a(u,v)+\mathfrak a(u+v,w)
 =\mathfrak a(v,w)+\mathfrak a(u,v+w).
\]
Likewise, comparing two parenthesisations of a triple product yields
\[
 \PhiMap(w)\mathfrak m(u,v)+\mathfrak m(uv,w)
 =\PhiMap(u)\mathfrak m(v,w)+\mathfrak m(u,vw).
\]
These identities are automatic, but they are valuable formalisation checks.

\subsection{The reverse sign}

For canonical ternary digits, base-$3$ normalisation replaces three units at
position $i$ by one unit at position $i+1$.  Evaluating at base $2$ makes that
operation smaller rather than larger.

\begin{theorem}[Reverse carry identities]\label{thm:reverse-carry}
For $u,v\in\Nzero$, let $\lambda_i$ be the carries in the relevant canonical
base-$3$ addition or multiplication.  Then
\begin{align*}
 \PsiMap(u+v)
   &=\PsiMap(u)+\PsiMap(v)-\sum_{i\ge0}\lambda_i2^i,\\
 \PsiMap(uv)
   &=\PsiMap(u)\PsiMap(v)-\sum_{i\ge0}\lambda_i2^i.
\end{align*}
Hence $\PsiMap$ is subadditive and submultiplicative.
\end{theorem}

\begin{proof}
The base-$3$ carry recurrence is
$r_i+\lambda_{i-1}=w_i+3\lambda_i$.  Evaluate at $2$ rather than $3$.
The shifted carry sum contributes $2\sum_i\lambda_i2^i$, leaving one negative
copy of the carry polynomial after rearrangement.
\end{proof}

\begin{warning}
Subadditivity and submultiplicativity do not imply monotonicity.  The explicit
counterexample $\PsiMap(8)=6>4=\PsiMap(9)$ remains essential.  Every later use
of $\PsiMap$ relies on its envelope or carry formula, never on an order claim.
\end{warning}

\section{Archimedean power envelopes}

The digit maps have exact growth exponents.  The proof is elementary but the
equality cases are the decisive feature.

\subsection{The general expanding-base theorem}

\begin{theorem}[Power envelope for radix reinterpretation]\label{thm:envelope-general}
Let $2\le b<c$ and put
\[
  \gamma=\frac{\log c}{\log b}>1.
\]
Then for every $n\in\N$,
\[
  \Rad_{b\to c}(n)\le n^\gamma.
\]
Equality holds if and only if $n=b^L$ for some $L\in\Nzero$.
\end{theorem}

\begin{proof}
Let $L=\floor{\log_b n}$ and write
\[
 Y_b=\frac{n}{b^L}
   =\sum_{j=0}^{L}d_{L-j}^{(b)}(n)b^{-j},
 \qquad
 Y_c=\frac{\Rad_{b\to c}(n)}{c^L}
   =\sum_{j=0}^{L}d_{L-j}^{(b)}(n)c^{-j}.
\]
The leading terms agree, and every lower place has $c^{-j}<b^{-j}$.
Therefore $Y_c\le Y_b$.  Since $Y_b\ge1$ and $\gamma>1$,
$Y_b\le Y_b^\gamma$.  As $b^{L\gamma}=c^L$,
\[
 \Rad_{b\to c}(n)=c^LY_c
 \le c^LY_b^\gamma=n^\gamma.
\]
For equality, first $Y_c=Y_b$, forcing every lower digit to vanish.  Then
$Y_b=Y_b^\gamma$, forcing the leading digit to be $1$.  Thus $n=b^L$.
The converse is immediate.
\end{proof}

\begin{theorem}[Reverse power envelope]\label{thm:reverse-envelope}
Under the same assumptions, put $\beta=1/\gamma$.  Then
\[
  \Rad_{c\to b}(m)\ge m^\beta
  \qquad(m\in\N),
\]
with equality if and only if $m=c^L$.
\end{theorem}

\begin{proof}
Normalise the canonical base-$c$ digits at their leading position:
\[
 Z_c=\frac{m}{c^L},\qquad
 Z_b=\frac{\Rad_{c\to b}(m)}{b^L}.
\]
Now $b^{-j}>c^{-j}$ at every lower place, so $Z_b\ge Z_c$.  Also
$Z_c\ge1\ge Z_c^{\beta-1}$, hence $Z_c\ge Z_c^\beta$.  Multiplication by
$b^L=c^{L\beta}$ proves the inequality.  Equality requires no lower digits
and leading digit $1$, exactly as before.
\end{proof}

Specialising gives the two fundamental bounds
\begin{equation}\label{eq:two-envelopes}
  \boxed{\PhiMap(n)\le n^{\thetaAE}},
  \qquad
  \boxed{\PsiMap(n)\ge n^{\alphaAE}}.
\end{equation}

\subsection{Exact decomposition of the forward gap}

The proof of \cref{thm:envelope-general} separates two independent sources of
strictness: curvature of $y\mapsto y^{\thetaAE}$ and contraction of lower
binary places when they are reweighted by powers of $3$.

\begin{proposition}[Curvature plus digit-weight gap]\label{prop:gap-decomposition}
Let $n\in\N$, $L=\floor{\log_2n}$, and
\[
  Y_2=\frac{n}{2^L},\qquad
  Y_3=\frac{\PhiMap(n)}{3^L}.
\]
Then
\begin{align}
 n^{\thetaAE}-\PhiMap(n)
 &=3^L\bigl[(Y_2^{\thetaAE}-Y_2)+(Y_2-Y_3)\bigr],
 \label{eq:gap-decomp}\\
 Y_2-Y_3
 &=\sum_{j=1}^{L}\varepsilon_{L-j}
     \bigl(2^{-j}-3^{-j}\bigr),
 \label{eq:digit-weight-gap}
\end{align}
where $\varepsilon_i$ are the binary digits of $n$.  In particular,
\begin{equation}\label{eq:gap-lower}
 n^{\thetaAE}-\PhiMap(n)
 \ge(\thetaAE-1)\left(\frac32\right)^L(n-2^L).
\end{equation}
\end{proposition}

\begin{proof}
The identity is obtained by adding and subtracting $3^LY_2$.  Formula
\eqref{eq:digit-weight-gap} is termwise.  Convexity gives
$Y_2^{\thetaAE}\ge1+\thetaAE(Y_2-1)$, hence
$Y_2^{\thetaAE}-Y_2\ge(\thetaAE-1)(Y_2-1)$, which becomes
\eqref{eq:gap-lower} after rescaling.
\end{proof}

The lower bound is especially strong near a power of $2$: if
$n=2^L+r$ with $1\le r<2^L$, then the gap is at least
$(\thetaAE-1)(3/2)^Lr$.  Thus the integer gap in a hypothetical solution is
not merely nonzero; it grows at a definite cross-radix scale.

\section{Specialisation to a hypothetical counterexample}

Assume from now through \cref{sec:why-not-closed} that
\begin{equation}\label{eq:candidate}
  M=2^x\in\N,
  \qquad A=3^x\in\N,
  \qquad x\notin\Z.
\end{equation}
Then $x>0$, $M>1$, $A>1$, and
\[
  A=M^{\thetaAE},\qquad M=A^{\alphaAE}.
\]
Neither $M$ is a power of $2$ nor $A$ a power of $3$.

\subsection{The four canonical sequences}

\begin{definition}[Cantor shadows and defects]\label{def:defects}
For $k\ge1$, set
\[
 P_k=\PhiMap(M^k),\qquad D_k=A^k-P_k,
\]
\[
 Q_k=\PsiMap(A^k),\qquad E_k=Q_k-M^k.
\]
\end{definition}

For bookkeeping at the identity element, we also set $P_0=Q_0=1$ and $D_0=E_0=0$.

\begin{proposition}[Strict integer sandwich]\label{prop:strict-sandwich}
For every $k\ge1$,
\[
  1\le D_k\le A^k-1,
  \qquad
  E_k\ge1.
\]
Equivalently,
\[
  \PhiMap(M^k)\le A^k-1,
  \qquad
  \PsiMap(A^k)\ge M^k+1.
\]
\end{proposition}

\begin{proof}
By \cref{thm:envelope-general},
$P_k\le(M^k)^{\thetaAE}=A^k$.  Equality would force $M^k$ to be a power of
$2$, hence $M$ a power of $2$, contrary to \eqref{eq:candidate}.  Both sides
are integers, so $D_k\ge1$.  The upper bound follows from $P_k\ge1$.
The reverse assertion follows from \cref{thm:reverse-envelope}; equality
would force $A^k$, and hence $A$, to be a power of $3$.
\end{proof}

\subsection{Synchronous digit length}

A small identity gives the radix construction its geometric coherence.

\begin{proposition}[Equal word length]\label{prop:same-length}
For every $k\ge1$, the binary word of $M^k$ and the ternary word of $A^k$
have the same length,
\[
  L_k+1,\qquad L_k=\floor{kx}.
\]
Moreover, $P_k$ has as its ternary word exactly the binary word of $M^k$.
Thus $P_k<A^k$ says that the ternary word of $A^k$ lies lexicographically
above the binary word of $M^k$ when the latter is read as a $0$--$1$ ternary
word.
\end{proposition}

\begin{proof}
Since $\log_2M=x$ and $\log_3A=x$,
\[
 \floor{\log_2M^k}=\floor{kx}=\floor{\log_3A^k}.
\]
The remaining statements are the definition of $\PhiMap$ and ordinary
positional comparison in base $3$.
\end{proof}

This lexical dominance is exact but, by itself, is not restrictive enough:
it is already implied by the universal power envelope.  Its value comes from
combining it with carries and valuations, not from treating it as a standalone
order miracle.

\subsection{Compatibility with integer translation}

Let $r\ge0$ and write $M=2^rM_0$, $A=3^rA_0$.  Then
\[
  \PhiMap(2^rn)=3^r\PhiMap(n),
  \qquad
  \PsiMap(3^rm)=2^r\PsiMap(m).
\]
Consequently
\[
 D_k(M,A)=3^{rk}D_k(M_0,A_0),
 \qquad
 E_k(M,A)=2^{rk}E_k(M_0,A_0).
\]
Thus the radix defects respect the baseline's integer-shift reduction exactly;
no information is lost by passing to a primitive translate.

\section{Exact defect dynamics and the carry budget}

\subsection{Forward recurrence}

Let $\carry_{k,\ell}$ denote the binary multiplication carry polynomial,
evaluated at $3$, for the product $M^kM^\ell=M^{k+\ell}$.  By
\cref{thm:mul-carry},
\[
  P_{k+\ell}=P_kP_\ell+\carry_{k,\ell},
  \qquad \carry_{k,\ell}\in\Nzero.
\]

\begin{theorem}[Forward defect recurrence]\label{thm:D-recurrence}
For $k,\ell\ge1$,
\begin{equation}\label{eq:D-recurrence}
 \boxed{
 D_{k+\ell}
 =A^kD_\ell+A^\ell D_k-D_kD_\ell-\carry_{k,\ell}.}
\end{equation}
In particular,
\begin{equation}\label{eq:D-square}
 D_{2k}=2A^kD_k-D_k^2-\carry_{k,k}.
\end{equation}
Because $D_{k+\ell}\ge1$, every hypothetical counterexample must satisfy the
carry-budget inequality
\begin{equation}\label{eq:carry-budget}
 \carry_{k,\ell}+1
 \le A^kD_\ell+A^\ell D_k-D_kD_\ell.
\end{equation}
\end{theorem}

\begin{proof}
Substitute $P_j=A^j-D_j$ into
$P_{k+\ell}=P_kP_\ell+\carry_{k,\ell}$ and expand.  The last inequality is
\cref{prop:strict-sandwich}.
\end{proof}

The recurrence has a useful interpretation.  If there were no binary carries,
then $P_{k+\ell}=P_kP_\ell$, and the defect would evolve by the first three
terms on the right of \eqref{eq:D-recurrence}.  Every actual carry consumes a
positive portion of that available defect.  A proof by this route must show
that, at some scale, the required carry expenditure exceeds the available
budget.

\subsection{Reverse recurrence}

Let $\revcarry_{k,\ell}$ be the base-$3$ multiplication carry polynomial,
evaluated at $2$, for $A^kA^\ell=A^{k+\ell}$.  Then
\[
 Q_{k+\ell}=Q_kQ_\ell-\revcarry_{k,\ell}.
\]

\begin{theorem}[Reverse defect recurrence]\label{thm:E-recurrence}
For $k,\ell\ge1$,
\begin{equation}\label{eq:E-recurrence}
 \boxed{
 E_{k+\ell}
 =M^kE_\ell+M^\ell E_k+E_kE_\ell-\revcarry_{k,\ell}.}
\end{equation}
Hence
\begin{equation}\label{eq:reverse-budget}
 \revcarry_{k,\ell}+1
 \le M^kE_\ell+M^\ell E_k+E_kE_\ell.
\end{equation}
\end{theorem}

\begin{proof}
Insert $Q_j=M^j+E_j$ into the reverse carry identity and use $E_{k+\ell}\ge1$.
\end{proof}

The signs in \eqref{eq:D-recurrence} and \eqref{eq:E-recurrence} are genuinely
different.  In the forward direction the mixed defect terms are reduced by
$D_kD_\ell$; in the reverse direction they are enlarged by $E_kE_\ell$.
This asymmetry is one reason to retain both sequences even though the forward
map has the cleaner valuation theory.

\subsection{Power bounds generated by the first defect}

Supermultiplicativity and submultiplicativity give immediate global bounds.

\begin{proposition}[Propagation from $k=1$]\label{prop:first-defect-prop}
For every $k\ge1$,
\begin{align*}
  (A-D_1)^k&\le P_k\le A^k-1,\\
  1\le D_k&\le A^k-(A-D_1)^k,\\
  M^k+1&\le Q_k\le(M+E_1)^k,\\
  1\le E_k&\le(M+E_1)^k-M^k.
\end{align*}
\end{proposition}

\begin{proof}
Repeated supermultiplicativity gives
$P_k=\PhiMap(M^k)\ge\PhiMap(M)^k=(A-D_1)^k$.
Repeated submultiplicativity gives
$Q_k\le\PsiMap(A)^k=(M+E_1)^k$.  The remaining bounds are definitions and
\cref{prop:strict-sandwich}.
\end{proof}

These estimates are not by themselves contradictory: the allowed intervals
still have exponential width.  Their main purpose is to calibrate any proposed
carry lower bound.

\section{Structural-prime laws}

\subsection{Valuation of a Cantor shadow}

Let
\[
  a=\vTwo(M),\qquad M=2^aU,\quad U\text{ odd}.
\]
Because the least nonzero binary digit of $M^k$ is in position $ak$, the least
nonzero ternary digit of $P_k$ is in the same position.

\begin{proposition}[Structural valuation]\label{prop:Pk-valuation}
For every $k\ge1$,
\[
  \vThree(P_k)=ak.
\]
If $U>1$, then
\begin{equation}\label{eq:Pk-anchor}
 \boxed{
 \vThree\bigl(P_k-3^{ak}\bigr)
 =ak+\vTwo(U^k-1).}
\end{equation}
The quotient after removal of this power of $3$ is congruent to $1$ modulo
$3$.
\end{proposition}

\begin{proof}
The shift identity gives $P_k=3^{ak}\PhiMap(U^k)$.  Since $U^k$ is odd, its
least binary digit is $1$, so $\PhiMap(U^k)$ is a $3$-adic unit.  If
$r=\vTwo(U^k-1)$, then the first binary digit after the least digit at which
$U^k$ differs from $1$ is position $r$, and that digit is $1$.  The isometry
\cref{thm:isometry}, or direct digit inspection, gives
$\vThree(\PhiMap(U^k)-1)=r$ with unit quotient $1$ modulo $3$.
\end{proof}

Let $b=\vThree(A)$.  Since $D_k=A^k-P_k$, unequal structural slopes force the
valuation of the defect without cancellation.

\begin{corollary}[Unequal-slope defect valuation]\label{cor:D-valuation}
If $a\ne b$, then for all $k\ge1$,
\[
  \vThree(D_k)=k\min(a,b).
\]
If $a=b$, then $3^{ak}\mid D_k$ and
\[
  \frac{D_k}{3^{ak}}
  =\left(\frac{A}{3^a}\right)^k-\PhiMap(U^k).
\]
\end{corollary}

\begin{proof}
When $a\ne b$, the two summands $A^k$ and $P_k$ have unequal $3$-adic
valuations $bk$ and $ak$, so their difference has the smaller valuation.  The
equal case follows by factoring $3^{ak}$.
\end{proof}

\subsection{The first collision in a binary square}

The next theorem identifies not just a lower bound but the exact first ternary
place occupied by the square-carry polynomial.

\begin{theorem}[Square first-collision theorem]\label{thm:first-collision}
Let $n=2^au$ with $u>1$ odd, and set $r=\vTwo(u-1)$.  Then
\[
  C(n):=\PhiMap(n^2)-\PhiMap(n)^2>0
\]
and
\begin{equation}\label{eq:first-collision}
  \boxed{
  \vThree(C(n))=2a+r,
  \qquad
  \frac{C(n)}{3^{2a+r}}\equiv1\pmod3.}
\end{equation}
\end{theorem}

\begin{proof}
The two lowest $1$-bits of $n$ occur in positions $a$ and $a+r$.  In the raw
binary convolution for $n^2$, the coefficient at position $2a$ is $1$, and
all coefficients through position $2a+r-1$ are at most $1$ with no incoming
carry.  At position $2a+r$, the two cross terms
$(a,a+r)$ and $(a+r,a)$ contribute exactly $2$.  Thus the first nonzero carry
leaves position $2a+r$, and that carry equals $1$.  The multiplication carry
identity evaluates the carry string in base $3$, proving
\eqref{eq:first-collision}.
\end{proof}

Apply this to $n=M^k=2^{ak}U^k$.

\begin{corollary}[Exact square-carry valuation along the orbit]
\label{cor:orbit-square-carry}
If $U>1$, then
\begin{equation}\label{eq:orbit-square-carry}
 \boxed{
 \vThree(P_{2k}-P_k^2)
 =2ak+\vTwo(U^k-1),}
\end{equation}
with leading $3$-adic unit congruent to $1$ modulo $3$.  Equivalently,
\begin{equation}\label{eq:defect-square-valuation}
 \boxed{
 \vThree\bigl(2A^kD_k-D_k^2-D_{2k}\bigr)
 =2ak+\vTwo(U^k-1).}
\end{equation}
Moreover,
\begin{equation}\label{eq:explicit-carry-budget}
 D_k(2A^k-D_k)
 \ge D_{2k}+3^{2ak+\vTwo(U^k-1)}
 \ge1+3^{2ak+\vTwo(U^k-1)}.
\end{equation}
\end{corollary}

\begin{proof}
Use \cref{thm:first-collision} and the identity
$P_{2k}-P_k^2=2A^kD_k-D_k^2-D_{2k}$.  Positivity and the least ternary place
give the real lower bound for the carry polynomial.
\end{proof}

For $k=2^j$, lifting the exponent gives an especially transparent law.  If
$U>1$ is odd and $j\ge1$, then
\[
 \vTwo(U^{2^j}-1)
 =\vTwo(U-1)+\vTwo(U+1)+j-1.
\]
Therefore
\begin{equation}\label{eq:dyadic-tower-valuation}
 \vThree(P_{2^{j+1}}-P_{2^j}^2)
 =2a\,2^j+j+\sigma_U,
 \qquad
 \sigma_U=\vTwo(U-1)+\vTwo(U+1)-1.
\end{equation}
This exact affine--exponential valuation profile is a compact target for Lean
formalisation and for any future $3$-adic interpolation argument.

\subsection{Pairwise valuation geometry and the resonance set}

The square theorem is one slice of a complete pairwise law.

\begin{theorem}[Pairwise Cantor-shadow distances]\label{thm:pairwise-shadow}
For distinct $k,\ell\ge0$,
\begin{equation}\label{eq:pairwise-shadow}
 \boxed{
 \vThree(P_k-P_\ell)=\vTwo(M^k-M^\ell).}
\end{equation}
\end{theorem}

\begin{proof}
This is \cref{thm:isometry} applied to $M^k$ and $M^\ell$.
\end{proof}

For $k>\ell$, define
\[
 R_{k,\ell}=\vTwo(M^k-M^\ell),
 \qquad
 S_{k,\ell}=\vThree(A^k-A^\ell).
\]
Since
\[
 P_k-P_\ell=(A^k-A^\ell)-(D_k-D_\ell),
\]
the ultrametric inequality determines the defect difference whenever the two
known valuations are unequal.

\begin{theorem}[Valuation resonance dichotomy]\label{thm:resonance}
If $R_{k,\ell}\ne S_{k,\ell}$, then
\begin{equation}\label{eq:defect-difference-valuation}
 \boxed{
 \vThree(D_k-D_\ell)=\min(R_{k,\ell},S_{k,\ell}).}
\end{equation}
If $S_{k,\ell}<R_{k,\ell}$, then additionally
\[
 D_k-D_\ell\equiv A^k-A^\ell\pmod{3^{R_{k,\ell}}}.
\]
All exceptional cancellation is therefore confined to the explicit resonance
locus
\[
  R_{k,\ell}=S_{k,\ell}.
\]
\end{theorem}

\begin{proof}
Set $X=A^k-A^\ell$ and $Y=D_k-D_\ell$.  We know
$\vThree(X-Y)=R_{k,\ell}$ and $\vThree(X)=S_{k,\ell}$.  If these valuations
are unequal, the non-Archimedean triangle law forces
$\vThree(Y)$ to be their minimum.  When $S<R$, the relation
$Y=X-(X-Y)$ also gives the displayed congruence.
\end{proof}

The resonance locus can be written explicitly by LTE.  Let $d=k-\ell$.
If $a=\vTwo(M)>0$, then
\[
  R_{k,\ell}=a\ell,
\]
because $M^d-1$ is odd.  If $a=0$, so $M$ is odd, then
\[
 R_{k,\ell}=
 \begin{cases}
  \vTwo(M-1),&d\text{ odd},\\
  \vTwo(M-1)+\vTwo(M+1)+\vTwo(d)-1,&d\text{ even}.
 \end{cases}
\]
Similarly, if $b=\vThree(A)>0$, then $S_{k,\ell}=b\ell$.  If $b=0$, then
\[
 S_{k,\ell}=
 \begin{cases}
  \vThree(A-1)+\vThree(d),&A\equiv1\pmod3,\\
  0,&A\equiv-1\pmod3\text{ and }d\text{ odd},\\
  \vThree(A^2-1)+\vThree(d),&A\equiv-1\pmod3\text{ and }d\text{ even}.
 \end{cases}
\]
Thus resonance is not an amorphous exceptional set: it lies on finitely many
linear relations among $\vTwo(d)$, $\vThree(d)$, and $\ell$.

\section{The mantissa observable and a sharp rotation gauge}

This section links the real size of $D_k$ to the irrational rotation
$k\mapsto\fracpart{kx}$.  Unlike a generic asymptotic estimate, the initial
formula is exact.

\subsection{A finite Cantor embedding of a dyadic mantissa}

Put
\[
  L_k=\floor{kx},\qquad t_k=\fracpart{kx}\in(0,1).
\]
Then
\[
  M^k=2^{L_k+t_k}=2^{L_k}(1+u_k),
  \qquad u_k=2^{t_k}-1.
\]
Because $M^k$ is an integer, $u_k$ is a dyadic rational with a terminating
binary expansion of length at most $L_k$:
\[
  u_k=\sum_{j=1}^{L_k}\varepsilon_{k,j}2^{-j}.
\]
Define its finite Cantor embedding
\[
  \Cantor(u_k)=\sum_{j=1}^{L_k}\varepsilon_{k,j}3^{-j}.
\]
The normalised Cantor shadow is then
\[
  \frac{P_k}{3^{L_k}}=1+\Cantor(u_k),
\]
whereas
\[
  \frac{A^k}{3^{L_k}}=3^{t_k}.
\]

\begin{theorem}[Exact mantissa formula]\label{thm:mantissa}
For every $k\ge1$,
\begin{equation}\label{eq:mantissa-ratio}
 \boxed{
 \frac{P_k}{A^k}
 =3^{-t_k}\bigl(1+\Cantor(2^{t_k}-1)\bigr),}
\end{equation}
and hence
\begin{equation}\label{eq:defect-observable}
 \boxed{
 \frac{D_k}{A^k}
 =1-3^{-t_k}\bigl(1+\Cantor(2^{t_k}-1)\bigr).}
\end{equation}
\end{theorem}

\begin{proof}
The binary digits below the leading bit of $M^k$ are exactly the terminating
binary digits of $u_k$.  Reinterpreting them in base $3$ gives
$P_k=3^{L_k}(1+\Cantor(u_k))$.  Divide by
$A^k=3^{L_k+t_k}$.
\end{proof}

\subsection{Universal two-sided bounds}

Termwise, $3^{-j}\le2^{-j}$, so
\[
  0\le\Cantor(u_k)\le u_k=2^{t_k}-1.
\]
This loses the fine digit structure but produces sharp, uniform control.

\begin{theorem}[Rotation gauge for the forward defect]\label{thm:rotation-gauge}
For every $k\ge1$,
\begin{equation}\label{eq:ratio-bounds}
  3^{-t_k}
  \le\frac{P_k}{A^k}
  \le\left(\frac23\right)^{t_k},
\end{equation}
so
\begin{equation}\label{eq:defect-bounds-exponential}
  1-\left(\frac23\right)^{t_k}
  \le\frac{D_k}{A^k}
  \le1-3^{-t_k}.
\end{equation}
In particular,
\begin{equation}\label{eq:defect-bounds-linear}
 \boxed{
  \frac{t_k}{3}
  \le\frac{D_k}{A^k}
  \le(\log3)t_k.}
\end{equation}
\end{theorem}

\begin{proof}
The bounds in \eqref{eq:ratio-bounds} follow from
$1\le1+\Cantor(u_k)\le1+u_k=2^{t_k}$.  Subtract from $1$ to obtain
\eqref{eq:defect-bounds-exponential}.  The function
$f(t)=1-(2/3)^t$ is concave on $[0,1]$, with $f(0)=0$ and $f(1)=1/3$, so
$f(t)\ge t/3$.  Finally $1-e^{-y}\le y$ with $y=t\log3$ gives the upper
linear bound.
\end{proof}

Equation \eqref{eq:defect-bounds-linear} is one of the central new reductions.
It says that the same integer $D_k$ is simultaneously

\begin{itemize}
\item an exact positive carry-budget variable in \eqref{eq:D-recurrence};
\item a $3$-adically constrained difference in
\cref{thm:resonance,cor:orbit-square-carry}; and
\item an Archimedean gauge of the rotation distance $t_k=\fracpart{kx}$,
within absolute constants independent of the hypothetical pair.
\end{itemize}

\subsection{Density and oscillation}

Because $x$ is irrational, the sequence $t_k=\fracpart{kx}$ is dense in
$[0,1]$.  The gauge immediately gives a nontrivial oscillation law.

\begin{corollary}[Macroscopic and thin relative defects]\label{cor:defect-oscillation}
For a hypothetical counterexample,
\[
  \liminf_{k\to\infty}\frac{D_k}{A^k}=0,
\]
and
\[
  \frac13\le
  \limsup_{k\to\infty}\frac{D_k}{A^k}
  \le\frac23.
\]
\end{corollary}

\begin{proof}
Choose a subsequence with $t_k\to0$ and use the upper bound in
\eqref{eq:defect-bounds-linear}.  Choose another with $t_k\to1$ and use the
lower bound in \eqref{eq:defect-bounds-exponential}.  The universal upper
bound follows from $1-3^{-t}<2/3$ for $t<1$.
\end{proof}

Thus $D_k$ is never a perturbative error sequence in absolute size.  It has
subsequences that are small relative to $A^k$, but it also repeatedly occupies
a fixed positive fraction of $A^k$.

\subsection{The reverse mantissa}

For completeness, write
\[
 A^k=3^{L_k}(1+v_k),\qquad v_k=3^{t_k}-1,
\]
and let
\[
  \mathcal G(v_k)=\sum_{j=1}^{L_k}\eta_{k,j}2^{-j}
\]
be the canonical ternary fractional digits of $v_k$, evaluated at binary
fractional place values.  Then
\[
  \frac{Q_k}{M^k}
  =2^{-t_k}\bigl(1+\mathcal G(3^{t_k}-1)\bigr).
\]
Since $2^{-j}\ge3^{-j}$ termwise,
$\mathcal G(v_k)\ge v_k$, and therefore
\begin{equation}\label{eq:E-lower}
  \frac{E_k}{M^k}
  \ge\left(\frac32\right)^{t_k}-1
  \ge(\log(3/2))t_k.
\end{equation}
An elementary first-nonzero-digit estimate also gives
\[
  \mathcal G(v)\le4v^{\alphaAE}\qquad(0<v<2),
\]
so, using $3^t-1\le2t$ on $[0,1]$,
\begin{equation}\label{eq:E-upper}
  \frac{E_k}{M^k}
  \le4(2t_k)^{\alphaAE}.
\end{equation}
The reverse defect is therefore controlled by a Hölder, rather than a linear,
window near $t=0$.  This loss reflects the digit $2$ and the nonmonotonicity of
$\PsiMap$.

\section{Diophantine consequences of the rotation gauge}

\subsection{An elementary absolute lower bound}

Let $L=L_k$.  Since $M^k$ is an integer strictly larger than $2^L$,
\[
  2^{t_k}-1=\frac{M^k-2^L}{2^L}\ge2^{-L}.
\]
Convexity of $2^t$ on $[0,1]$ gives $2^t\le1+t$, hence $t_k\ge2^{-L}$.
Combining with \eqref{eq:defect-bounds-linear} yields
\[
  D_k\ge\frac{A^k}{3\,2^{L_k}}
       \ge\frac13\left(\frac32\right)^{L_k}.
\]
The sharper digit-sensitive bound \eqref{eq:gap-lower} is
\[
 D_k\ge(\thetaAE-1)
 \left(\frac32\right)^{L_k}\bigl(M^k-2^{L_k}\bigr).
\]
Both are unconditional consequences of the radix envelope once the output
$A^k$ is integral.

\subsection{Baker--Matveev amplification}

The elementary bound allows $t_k$ to be exponentially small in $k$.  For a
fixed integer $M$ not equal to a power of $2$, linear forms in logarithms rule
out such extreme approximation.

\begin{theorem}[Effective polynomial lower bound]\label{thm:baker-defect}
Fix an integer $M>1$ that is not a power of $2$, and put $x=\log_2M$.
There exist effectively computable constants $c(M)>0$ and $C(M)>0$ such that
for every $k\ge1$,
\[
  \fracpart{kx}\ge c(M)k^{-C(M)}.
\]
Consequently, if $A=M^{\thetaAE}$ is an integer, then
\begin{equation}\label{eq:baker-D}
  D_k\ge\frac{c(M)}{3}A^kk^{-C(M)}.
\end{equation}
\end{theorem}

\begin{proof}
Let $L_k=\floor{kx}$.  The nonzero linear form
\[
 \Lambda_k=k\log M-L_k\log2=(\log2)\fracpart{kx}
\]
has fixed algebraic logarithms and integer coefficients of size $O(k)$.
A standard explicit lower bound of Baker--W\"ustholz or Matveev gives
$|\Lambda_k|\ge c'(M)k^{-C(M)}$ after adjusting constants for finitely many
small $k$.  Apply \eqref{eq:defect-bounds-linear}.
\end{proof}

The point is qualitative but important: on every scale, the positive integer
$D_k$ is exponentially large up to a fixed polynomial factor.  Therefore a
closing argument cannot treat $D_k$ as a small absolute error.  It must use
its special digit or $3$-adic structure.

\subsection{Infinitely many thin relative defects}

Continued-fraction convergents to an irrational number alternate around it.
There are therefore infinitely many $k$ with
\[
  0<\fracpart{kx}<\frac1k.
\]
For those $k$, the upper gauge gives
\begin{equation}\label{eq:thin-subsequence}
  1\le D_k<\frac{\log3}{k}A^k.
\end{equation}
Thus the same sequence has Baker-controlled lower size at every index and
Dirichlet-thin relative size on an infinite subsequence.  Any successful
cross-place argument should be tested first on this subsequence.

\section{A nonlinear spectral-radius reformulation}

The carry inequalities naturally suggest a spectral point of view.

\begin{definition}[Radix spectral radii]
For $n,m\ge1$, define
\[
 r_\Phi(n)=\lim_{k\to\infty}\PhiMap(n^k)^{1/k},
 \qquad
 r_\Psi(m)=\lim_{k\to\infty}\PsiMap(m^k)^{1/k},
\]
provided the limits exist.
\end{definition}

\begin{theorem}[Exact radix spectral radii]\label{thm:spectral-radius}
For every $n,m\in\N$,
\begin{align}
 r_\Phi(n)&=n^{\thetaAE}
 =\sup_{k\ge1}\PhiMap(n^k)^{1/k},
 \label{eq:phi-spectral}\\
 r_\Psi(m)&=m^{\alphaAE}
 =\inf_{k\ge1}\PsiMap(m^k)^{1/k}.
 \label{eq:psi-spectral}
\end{align}
\end{theorem}

\begin{proof}
The upper envelope gives
$\PhiMap(n^k)\le n^{k\thetaAE}$.  If $L=\floor{\log_2n^k}$, then
$\PhiMap(n^k)\ge3^L>n^{k\thetaAE}/3$.  Taking $k$th roots squeezes the limit
to $n^{\thetaAE}$.  Supermultiplicativity of $\PhiMap$ makes
$\log\PhiMap(n^k)$ superadditive in $k$, so Fekete's lemma identifies the
limit with the supremum of the normalised terms.

For $\PsiMap$, the reverse envelope gives the lower bound
$m^{k\alphaAE}$.  If the ternary length is $L+1$, then
\[
 \PsiMap(m^k)\le2\sum_{i=0}^{L}2^i<2^{L+2}<4m^{k\alphaAE}.
\]
The $k$th roots again squeeze the limit.  Submultiplicativity and Fekete's
lemma for subadditive logarithms give the infimum formula.
\end{proof}

\begin{corollary}[Spectral-radius form of Alaoglu--Erd\H{o}s]
\label{cor:spectral-AE}
The Alaoglu--Erd\H{o}s conjecture is equivalent to
\[
  \boxed{
  r_\Phi(M)\in\N\quad\Longrightarrow\quad M=2^j
  \text{ for some }j\in\Nzero.}
\]
Equivalently: the multiplicative orbit of an integer has an integer
binary-to-ternary radix spectral radius only when every power belongs to the
carry-free one-bit orbit.
\end{corollary}

\begin{proof}
By \cref{thm:spectral-radius}, $r_\Phi(M)=M^{\thetaAE}$.  If this is the
integer $A$, then with $x=\log_2M$ we have $2^x=M$ and
$3^x=M^{\thetaAE}=A$.  The original conjecture says that $x$ is integral,
which is equivalent to $M$ being a power of $2$.
\end{proof}

This reformulation is not presented as a proof by renaming.  Its value is that
$P_k=\PhiMap(M^k)$ is generated by a fixed multiplication-by-$M$ digit
transducer, while $r_\Phi(M)$ is its orbit growth rate under the Archimedean
weight $3^i$.  It opens a route through transfer operators, pressure, or
infinite-state Perron--Frobenius theory that is absent from a logarithm-only
formulation.

\section{The adelic shadow orbit}

The construction can be summarised as a single orbit viewed at two places.

\subsection{Two compatible geometries}

The binary orbit
\[
  1,M,M^2,M^3,\ldots\subset\Z_2
\]
has the Cantor shadow
\[
  1,P_1,P_2,P_3,\ldots\subset\Z_3,
  \qquad P_k=\widehat\Phi(M^k).
\]
By \cref{thm:pairwise-shadow}, all pairwise distances are transported exactly:
\[
  |P_k-P_\ell|_3=3^{-\vTwo(M^k-M^\ell)}.
\]
At the Archimedean place the shadow sits strictly below the geometric orbit
$A^k$:
\[
  0<P_k<A^k,
  \qquad A^k-P_k=D_k.
\]
The same $D_k$ is then governed by the rotation gauge
\eqref{eq:defect-bounds-linear}.

This produces an exact three-way dictionary:
\[
\begin{array}{ccc}
 M^k\in\Z_2
 &\xrightarrow{\ \widehat\Phi\ }
 &P_k\in\Cantor_3\subset\Z_3\\[0.4em]
 &&\Big\downarrow\scriptstyle +D_k\\[-0.1em]
 &&A^k\in\Z\subset\R\cap\Z_3.
\end{array}
\]
The horizontal arrow is an ultrametric isometry; the vertical arrow is a
positive Archimedean displacement with exact nonlinear recurrences.

\subsection{Why this is more than the compact-orbit reformulation}

The baseline already records compact and $S$-integral orbit formulations.
The radix shadow adds two pieces of data not supplied by compactness alone:

\begin{enumerate}
\item the point $P_k$ is not merely close to an orbit; it belongs to a rigid
Cantor subset and carries the complete binary digit word of $M^k$;
\item the valuation of every difference $P_k-P_\ell$ is known exactly, not
just bounded.
\end{enumerate}

This exactness is what makes \cref{thm:resonance,thm:first-collision}
possible.  The unresolved task is to make the Archimedean order
$P_k<A^k$ interact decisively with the ultrametric isometry.

\section{Audited dead ends and the precise missing step}\label{sec:why-not-closed}

The radix mechanism is strong enough to generate seductive false proofs.  It
is useful to record exactly where they fail.

\subsection{The envelope alone is universal}

The inequality $\PhiMap(M^k)\le A^k$ may look like a new restriction on a
counterexample.  It is not: $\PhiMap(n)\le n^{\thetaAE}$ holds for every
integer $n$.  The genuinely arithmetic input is that $A^k$ is an integer, so
$D_k$ is an integer, and that all $D_k$ belong to one coupled carry recurrence.
Any argument using only the pointwise envelope is a reformulation, not a
closure.

\subsection{The reverse map is not monotone}

From $P_k<A^k$ one cannot apply $\PsiMap$ and retain the inequality.  The
example $8<9$ but $\PsiMap(8)>\PsiMap(9)$ is the smallest visible obstruction.
The valid reverse inequality $\PsiMap(A^k)>M^k$ comes independently from the
reverse power envelope.

\subsection{Positive carries can absorb the entire discrepancy}

Equation \eqref{eq:D-recurrence} does not make $D_{k+\ell}$ grow
monotonically.  A large carry polynomial subtracts from the nominal defect
budget.  Conversely, the fact that the carry is large is not itself a
contradiction: its real size can be exponential while its first nonzero
ternary place is only logarithmic in $k$.

\subsection{Large $3$-adic divisibility does not imply real smallness}

The exact valuation
\[
 \vThree(P_{2k}-P_k^2)=2ak+\vTwo(U^k-1)
\]
can grow without forcing the positive integer $P_{2k}-P_k^2$ to vanish.  A
product-formula argument would need an independent upper bound below the
corresponding power of $3$.  The current Archimedean bounds are far too large
for that.

\subsection{The two defects do not automatically descend}

Both $D_k<A^k$ and, from a crude digit bound, $E_k<2M^k$ for large scales.
This suggests a descent from $(M^k,A^k)$ to $(E_k,D_k)$.  But no identity
$D_k=E_k^{\thetaAE}$ is available, and the carry corrections prevent the two
radix maps from forming an order adjunction on all integers.  A descent claim
without a preserved exponential relation is invalid.

\subsection{A precise closing lemma}

The new material reduces the desired contradiction to a cross-place rigidity
statement of the following kind.

\begin{problem}[Carry--rotation rigidity]\label{prob:closing}
Classify integer pairs $M,A>1$ for which there exist positive integer
sequences $D_k$ and nonnegative carry polynomials $\carry_{k,\ell}$ satisfying
simultaneously
\begin{enumerate}[label=(\roman*)]
\item $A=M^{\thetaAE}$;
\item the exact recurrence \eqref{eq:D-recurrence};
\item the rotation gauge \eqref{eq:defect-bounds-linear};
\item the cross-prime isometry constraints
\eqref{eq:pairwise-shadow}--\eqref{eq:defect-difference-valuation}; and
\item the first-collision law \eqref{eq:defect-square-valuation} with leading
unit $1$.
\end{enumerate}
Prove that the only possibilities have $M=2^j$ and $A=3^j$.
\end{problem}

Condition (i) is still the original transcendental relation, so
\cref{prob:closing} is not advertised as a logically weaker theorem.  The
point is methodological: conditions (ii)--(v) expose rigid, computable
consequences that a proof may attack by $p$-adic interpolation, transfer
operators, or carry combinatorics.  The baseline formulations do not package
these consequences in one sequence.

\section{Promising next attacks}

\subsection{A transfer operator for multiplication by $M$}

Multiplication by a fixed $M$ is a finite-carry transducer when binary digits
are read from least significant to most significant.  Iterating it generates
the words of $M^k$.  The functional
\[
  w\longmapsto\sum_i w_i3^i
\]
assigns an Archimedean weight to the output word.  The exact spectral radius
\cref{thm:spectral-radius} suggests constructing an infinite-dimensional
positive transfer operator whose orbit vector is $P_k$.  The desired rigidity
would follow from a theorem of the schematic form:

\begin{conjecture}[Integer pressure rigidity]\label{conj:pressure}
If the orbit pressure of the multiplication-by-$M$ binary carry transducer,
under ternary place weights, is the logarithm of an integer, then the
transducer orbit from $1$ is eventually carry-free.  Consequently $M$ is a
power of $2$.
\end{conjecture}

A finite matrix representation is not presently justified; claiming one
would be false in general because digit sums of powers need not satisfy a
fixed linear recurrence.  The plausible object is a transfer operator on a
space of carry-state distributions or $2$-adic cylinders.

\subsection{Exploit the sparse resonance locus}

Outside $R_{k,\ell}=S_{k,\ell}$, the valuation of $D_k-D_\ell$ is exact.  The
LTE formula shows that resonance lies on explicit diagonals in
$(\vTwo(k-\ell),\vThree(k-\ell),\ell)$.  A possible route is:

\begin{enumerate}
\item choose many indices with prescribed $2$- and $3$-adic valuations but
with fractional parts $\fracpart{kx}$ in a short real interval;
\item use \eqref{eq:defect-bounds-linear} to confine the corresponding
$D_k$ to short Archimedean intervals;
\item use \cref{thm:resonance} to force their pairwise differences into
incompatible $3$-adic residue classes.
\end{enumerate}

The difficulty is quantitative.  Ordinary equidistribution gives real
intervals of polynomial resolution, whereas forcing a positive integer to
vanish from $3$-adic divisibility needs exponential resolution.  A successful
version must gain cancellation by combining many differences, not by a single
pair.

\subsection{Carry-polynomial $S$-unit estimates}

The square carry
\[
  C_k=P_{2k}-P_k^2
\]
has a ternary expansion whose first nonzero digit is known exactly.  Its
remaining digits are generated by a constrained convolution of the binary
support of $M^k$.  One can ask whether infinitely many unusually sparse carry
polynomials are possible when $M$ has an odd factor.  Results on digit sums of
powers and $S$-unit equations suggest the following intermediate target.

\begin{problem}[Sparse-carry exclusion]
Prove an effective lower bound for the number or weighted span of nonzero
ternary digits of $C_k$ for $M$ not a power of $2$, strong enough to conflict
with the carry budget \eqref{eq:carry-budget} on a Dirichlet-thin subsequence.
\end{problem}

A mere statement that the digit sum tends to infinity is too weak; the budget
is exponential.  The needed estimate must retain the position and
multiplicity of carries.

\subsection{A two-place determinant}

The baseline's determinant methods are Archimedean.  The isometry
\cref{thm:isometry} makes it natural to build determinants from several
$P_k=A^k-D_k$ such that

\begin{itemize}
\item row or column differences acquire prescribed high powers of $3$ from
the binary orbit; while
\item the Archimedean size is reduced by choosing $k$ with clustered
fractional parts $\fracpart{kx}$.
\end{itemize}

The target is a nonzero integer determinant whose real absolute value is less
than $1$, with its nonvanishing certified by the first differing Cantor digit.
This is more specific than a generic Vandermonde repetition: the $3$-adic
valuation matrix is known exactly from $M^k-M^\ell$.

\subsection{Certified computation as theorem discovery}

A useful search should track more than near-integrality of $M^{\thetaAE}$.
For each candidate integer $M$ and its only possible nearby $A$, compute
\[
  D_1,D_2,D_4,\ldots,
\]
the exact carry polynomials, their first-collision positions, and the
resonance matrix $R_{k,\ell}-S_{k,\ell}$.  The goal is to discover a finite
inequality or congruence pattern that can later be proved symbolically.  This
is complementary to the baseline's directed-rounding scan and should not be
reported as a larger zero-free theorem until certificates are kernel-replayed.

\section{Lean formalisation plan}

The results above are designed to separate cleanly into small modules.  The
following dependency order avoids importing the large existing corpus until
the final candidate specialisation.

\subsection{Module 1: canonical digit evaluation}

Suggested definitions:

\begin{lstlisting}[language={}] 
-- schematic Lean, not claimed to compile verbatim

def evalDigits (base : Nat) (d : List Nat) : Nat :=
  d.foldr (fun x acc => x + base * acc) 0

def reinterpret (source target n : Nat) : Nat :=
  evalDigits target (Nat.digits source n)

def phi : Nat -> Nat := reinterpret 2 3

def psi : Nat -> Nat := reinterpret 3 2
\end{lstlisting}

Core theorems are canonical-digit reconstruction, shift compatibility, and
strict monotonicity of \texttt{reinterpret b c} under $b<c$.

\subsection{Module 2: valuation isometry}

Use integers for subtraction:

\begin{lstlisting}[language={}]
theorem padicVal_phi_sub
    {m n : Nat} (h : m != n) :
    padicValInt 3 ((phi m : Int) - phi n)
      = padicValInt 2 ((m : Int) - n)
\end{lstlisting}

The proof should avoid general $p$-adic analysis: identify the least differing
digit and factor the difference as $3^r$ times a unit.

\subsection{Module 3: carry normalisation}

Define a generic normaliser for a finite coefficient list in base $b$,
returning both output digits and carries.  Prove a polynomial-evaluation
identity once:
\[
 \operatorname{eval}_c(\operatorname{normalise}_b r)
 =\operatorname{eval}_c(r)+(c-b)\operatorname{eval}_c(\operatorname{carry}).
\]
Addition and multiplication are then corollaries.  This general statement is
less brittle than separate bit-level proofs.

\subsection{Module 4: real envelopes}

The only analytic inputs are $\log$, real powers, and convexity of
$y\mapsto y^\gamma$ on $[1,\infty)$.  A useful intermediate lemma is the
normalised digit comparison
\[
 \frac{\Rad_{b\to c}(n)}{c^L}\le\frac{n}{b^L}.
\]
Equality classification should be formalised at the same time; downstream
strict positivity of $D_k$ depends on it.

\subsection{Module 5: candidate defects}

Import the existing candidate predicate only here.  Define $P_k,D_k,Q_k,E_k$
and prove the recurrences by ring normalisation.  The carry-budget inequalities
then require only positivity of the strict defects.

\subsection{Module 6: first collision and LTE}

Formalise the two lowest set-bit positions of $2^au$.  The key local statement
is that all square-convolution coefficients before $2a+r$ normalise without a
carry and the coefficient at $2a+r$ is exactly $2$.  Combine this with
Mathlib's LTE lemmas to obtain \eqref{eq:dyadic-tower-valuation}.

\subsection{Module 7: mantissa gauge}

Avoid formalising a global discontinuous Cantor function initially.  For each
$k$, use the finite digit list of $M^k$ and define its normalised binary and
ternary evaluations.  The inequalities
\[
  1\le Y_3\le Y_2=2^{\fracpart{kx}}
\]
are sufficient for \eqref{eq:defect-bounds-linear}.  The exact Cantor formula
can be added as an interface theorem after the finite proof is stable.

\subsection{Proposed theorem inventory}

A compact first milestone is:

\begin{longtable}{p{0.39\textwidth}p{0.51\textwidth}}
\toprule
Suggested Lean name & Mathematical statement\\
\midrule
\endhead
\texttt{phi\_strictMono} & \cref{prop:order} for bases $2,3$.\\
\texttt{padicVal\_phi\_sub\_eq} & \cref{thm:isometry}.\\
\texttt{phi\_add\_eq\_add\_carry} & \cref{thm:add-carry}.\\
\texttt{phi\_mul\_eq\_mul\_carry} & \cref{thm:mul-carry}.\\
\texttt{psi\_add\_eq\_sub\_carry} & addition part of \cref{thm:reverse-carry}.\\
\texttt{psi\_mul\_eq\_sub\_carry} & multiplication part of \cref{thm:reverse-carry}.\\
\texttt{phi\_le\_rpow} & $\PhiMap(n)\le n^{\thetaAE}$.\\
\texttt{phi\_eq\_rpow\_iff} & equality iff $n$ is a power of $2$.\\
\texttt{psi\_ge\_rpow} & $\PsiMap(n)\ge n^{\alphaAE}$.\\
\texttt{candidate\_defect\_pos} & \cref{prop:strict-sandwich}.\\
\texttt{defect\_mul\_recurrence} & \eqref{eq:D-recurrence}.\\
\texttt{squareCarry\_padicVal} & \cref{thm:first-collision}.\\
\texttt{candidate\_defect\_rotation\_bounds} & \eqref{eq:defect-bounds-linear}.\\
\texttt{radixSpectralRadius} & \cref{thm:spectral-radius}.\\
\texttt{alaogluErdos\_iff\_integer\_radixRadius} & \cref{cor:spectral-AE}.\\
\bottomrule
\end{longtable}

None of these identifiers is claimed to exist in the current repository.  The
list is a formalisation specification, not a status report.

\section{Independent executable audit}

A short Python program was used only to detect indexing and sign mistakes.
The exact proofs above do not depend on it.  The audit performed the following
checks:

\begin{itemize}
\item \cref{thm:isometry} for every pair $0\le n<m<2048$;
\item both addition and multiplication carry identities, in both radix
directions, for every operand pair $0\le u,v<384$;
\item \cref{thm:first-collision} for every $1<n<250000$ that is not a power
of $2$;
\item the two real envelopes for every $1\le n\le200000$, using
$100$-decimal-digit arithmetic as a numerical sanity check.
\end{itemize}

All checks passed. \evidence{C}  The central routines are reproduced below so
that the sign convention is auditable.

\begin{lstlisting}[language=Python]
def digits(n, base):
    if n == 0:
        return [0]
    out = []
    while n:
        n, r = divmod(n, base)
        out.append(r)          # least significant first
    return out

def eval_digits(ds, base):
    acc = 0
    for d in reversed(ds):
        acc = acc * base + d
    return acc

def reinterpret(n, source, target):
    return eval_digits(digits(n, source), target)

def mul_carry_poly(u, v, source, target):
    du, dv = digits(u, source), digits(v, source)
    raw = [0] * (len(du) + len(dv) + 4)
    for i, a in enumerate(du):
        for j, b in enumerate(dv):
            raw[i + j] += a * b
    carry = 0
    poly = 0
    for i in range(len(raw)):
        total = raw[i] + carry
        carry = total // source
        poly += carry * target**i
    return poly

phi = lambda n: reinterpret(n, 2, 3)
psi = lambda n: reinterpret(n, 3, 2)

# Forward sign: positive carry polynomial.
assert phi(u*v) == phi(u)*phi(v) + mul_carry_poly(u,v,2,3)
# Reverse sign: negative carry polynomial.
assert psi(u*v) == psi(u)*psi(v) - mul_carry_poly(u,v,3,2)
\end{lstlisting}

For a kernel-verified development, the finite audit should be replaced by the
module plan of the previous section, not imported through \texttt{native\_decide} as a
substitute for the general proofs.

\section{Worked symbolic examples}

\subsection{A single binary carry}

Take $u=v=1$.  Binary addition gives $1+1=10_2$.  Hence
\[
 \PhiMap(2)=3,
 \qquad \PhiMap(1)+\PhiMap(1)=2,
\]
and the carry polynomial is $1$.  The difference $3-2$ is exactly the gain
from replacing two units of $3^0$ by one unit of $3^1$.

\subsection{A carry-free and a carrying square}

For $n=5=(101)_2$,
\[
 \PhiMap(5)=10,
 \qquad 25=(11001)_2,
 \qquad \PhiMap(25)=109.
\]
Thus
\[
 C(5)=109-10^2=9=3^2.
\]
Here $\vTwo(5-1)=2$, agreeing with \cref{thm:first-collision}.

For $n=2^a$, every power has one binary digit, so
\[
 \PhiMap(n^2)=\PhiMap(n)^2
\]
and the carry polynomial vanishes.  The equality case of the power envelope
identifies precisely these monomial orbits.

\subsection{Why the reverse map can fall}

The canonical ternary increment
\[
 (22)_3+1=(100)_3
\]
contains a two-position carry cascade.  Evaluating before normalisation at
base $2$ gives $6+1=7$; evaluating after normalisation gives $4$.  The reverse
carry polynomial is $3$, exactly accounting for the loss.  This is the local
mechanism behind $\PsiMap(8)=6>4=\PsiMap(9)$.

\section{Conclusions}

This continuation does not produce a completed proof of the
Alaoglu--Erd\H{o}s conjecture.  It does produce a coherent new rigidity package
that is independent of the baseline's main routes:

\begin{enumerate}
\item The binary-to-ternary digit map is simultaneously an Archimedean order
embedding and an exact $2$-adic/$3$-adic isometry.
\item Addition and multiplication have explicit positive carry cocycles; the
reverse digit evaluation has explicit negative cocycles.
\item The smooth curve $y=x^{\thetaAE}$ is an exact upper envelope of the
Cantor graph $y=\PhiMap(x)$, touching it only at powers of $2$.
\item A hypothetical counterexample generates positive integer defects
$D_k,E_k$ with coupled nonlinear recurrences.
\item The relative forward defect is comparable, with absolute constants, to
$\fracpart{kx}$, while its pairwise $3$-adic geometry is inherited exactly
from the binary power orbit.
\item The first square carry has a computable location and unit coefficient,
producing exact valuation and real carry-budget constraints.
\item The conjecture is equivalent to integer rigidity of a nonlinear radix
spectral radius.
\end{enumerate}

The most promising next move is not another isolated inequality.  It is a
multi-index construction that combines the exact $3$-adic distance matrix of
the Cantor shadows with clustered real rotation phases, so that many defect
differences cancel Archimedeanly while retaining certified nonzero
$3$-adic leading terms.  That is a concrete determinant/interpolation problem
and a natural target for both Lean formalisation and experimental theorem
discovery.

\begin{thebibliography}{99}

\bibitem{Unified}
Project manuscript,
\emph{Unified report on the Alaoglu--Erd\H{o}s conjecture},
attached consolidated report, integrated through continuation 27, 2026.

\bibitem{Baker}
A. Baker,
\emph{Transcendental Number Theory},
Cambridge University Press, 1975.

\bibitem{BakerWustholz}
A. Baker and G. W\"ustholz,
\emph{Logarithmic Forms and Diophantine Geometry},
New Mathematical Monographs 9, Cambridge University Press, 2007.

\bibitem{Matveev}
E. M. Matveev,
An explicit lower bound for a homogeneous rational linear form in logarithms
of algebraic numbers. II,
\emph{Izvestiya: Mathematics} 64 (2000), 1217--1269.

\bibitem{KuipersNiederreiter}
L. Kuipers and H. Niederreiter,
\emph{Uniform Distribution of Sequences},
Wiley, 1974; reprint, Dover, 2006.

\bibitem{AlloucheShallit}
J.-P. Allouche and J. Shallit,
\emph{Automatic Sequences: Theory, Applications, Generalizations},
Cambridge University Press, 2003.

\bibitem{SengeStraus}
H. G. Senge and E. G. Straus,
PV-numbers and sets of multiplicity,
\emph{Periodica Mathematica Hungarica} 3 (1973), 93--100.

\end{thebibliography}

\end{document}
